\chapter{总结与展望}
\section{论文工作总结}

本文围绕"高表现力三维数字人舞蹈内容生成"这一核心课题，从全身舞蹈动作生成与数字人外观生成编辑两个维度展开系统性研究，提出了一套兼顾音乐对齐舞蹈动作生成与角色外观快速定制的技术框架。论文的主要工作与贡献总结如下：

在三维数字人全身舞蹈动作生成方面，本文针对现有舞蹈数据集在动作精度、细节完整性与风格多样性方面的不足，采用专业的全身动作捕捉系统构建了高质量舞蹈动作数据集SoulDance。该数据集包含12.5小时成对音乐与舞蹈数据，在同一时间轴下同时提供身体动作、手部姿态与面部表情信息，覆盖284段音乐片段与15种音乐风格，由专业舞者在执行导演全程监督下完成录制，并以BVH骨架文件、SMPL-X人体模型与FLAME面部参数等多种表示形式发布，在数据规模、动作精度与细节完整性上实现了全面均衡。在此数据基础上，本文提出全身舞蹈生成框架SoulNet，其核心包含三个创新模块：分层残差向量量化编码（HRVQ）在残差向量量化的基础上构建"身体—手部—面部"的多层链式结构，通过跨部位交互显式建模三者之间的空间关系与层级依赖，有效提升了动作编码的精度与全身各部位的协调性；音乐对齐生成模型（MAGM）采用基于Transformer的掩码预测与残差细化两阶段架构，将层级化舞蹈单元组合为全局连贯的舞蹈序列；预训练的音乐动作检索模块（MMR）利用大规模公开音乐—舞蹈数据学习跨模态对齐先验，并以对比学习损失作为生成过程的对齐引导，从而在节拍结构与情绪张力上提升音乐与舞蹈的一致性。此外，本文还提出了情感对齐分数（EAS）与音乐动作匹配分数（MMR-MS）两项新评估指标，用于更精确地评估面部表情与音乐情绪、舞蹈动作与音乐节奏之间的对齐程度。在SoulDance与AIST++数据集上的实验结果表明，SoulNet在FID、多样性、音乐对齐等多项指标上均取得了当前最优性能，生成的全身舞蹈序列在自然性、协调性与表现力方面显著优于现有方法。

在三维数字人外观生成与编辑方面，本文针对现有外观个性化编辑方法在训练代价高、外观身份保持受限以及多视角一致性不足等方面的局限，提出了基于视角迭代自注意力控制的无训练个性化视觉编辑框架VisCtrl。该方法通过DDIM反演分别获取参考图像与目标图像的噪声轨迹，在去噪过程中采用迭代式自注意力机制，将参考主体的关键外观特征逐步注入目标图像，同时利用交叉注意力机制维持目标图像的整体结构不变。为将二维编辑能力拓展至三维数字人场景，本文进一步提出特征渐进采样（FGS）策略，通过对参考图像潜变量特征的随机采样与加权融合，在多视角之间实现一致的外观特征注入，有效缓解了跨视角编辑不一致导致的三维重建伪影问题。实验结果表明，VisCtrl在二维图像编辑、视频逐帧编辑以及基于3DGS的三维数字人外观编辑任务中，均在CLIP-I、BG LPIPS、SSIM等指标上取得了优于现有方法的表现，仅需单张参考图像即可完成高效的个性化外观定制，具备即插即用与模块化组合的优势。该方法为全身舞蹈动作生成提供了风格明确且视觉一致的角色外观基础，使数字人在舞蹈表演中能够兼具个性化外观与高表现力动作。

综上所述，本文从舞蹈动作生成与外观个性化编辑两个方向展开协同研究，构建了一条从音乐驱动全身舞蹈生成到角色外观快速定制的完整技术路径，为虚拟偶像表演、沉浸式社交与数字内容生产等应用场景提供了兼顾动作表现与视觉表达的技术支撑。

\section{未来工作展望}

尽管本文在三维数字人舞蹈内容生成方面取得了一定的研究成果，但受限于当前技术条件与研究周期，仍存在若干值得深入探索的方向。结合本文工作的局限性与领域发展趋势，未来研究可从以下几个方面展开：

首先，面向长序列与复杂编舞结构的舞蹈生成是提升舞蹈动作生成实用性的关键方向。当前SoulNet主要在中等长度的舞蹈片段上进行训练与评估，对于超长序列（如完整的三至五分钟舞蹈表演）的生成，仍面临动作质量随时长衰减、全局编舞结构缺失等挑战。未来可引入层次化的时序建模策略，例如在高层建模段落级的编舞结构与情绪弧线，在低层细化帧级的关节运动，从而在保证局部动作质量的同时维持全局的结构连贯性与叙事完整性。此外，结合大语言模型对编舞意图的理解能力，实现基于自然语言描述的舞蹈编排控制，也是一个具有应用前景的研究方向。

其次，外观编辑与舞蹈动作的端到端联合生成是实现三维数字人整体表现力统一的重要目标。本文目前将外观个性化编辑与舞蹈动作生成作为两个相对独立的模块分别进行研究，尚未实现二者在统一框架下的端到端联合优化。在实际应用中，角色的外观风格与舞蹈动作之间存在内在的语义关联——例如，特定服饰造型可能对应特定的舞蹈风格，角色体型特征也会影响动作的物理合理性与视觉表现。未来可探索将外观表示与动作表示统一到共享的潜在空间中，通过联合建模实现"外观—动作"的协同生成与风格一致性约束，从而进一步提升三维数字人整体表现的连贯性与艺术感染力。

最后，多模态条件的融合与可控生成是拓宽创作自由度的核心方向。本文主要以音乐作为舞蹈生成的条件输入，而在实际创作场景中，用户可能希望通过文本描述、情绪标签、风格参考视频等多种模态对生成结果进行联合控制。未来可在现有音乐条件对齐机制的基础上，引入多模态融合策略，支持文本—音乐—风格的多条件联合驱动，使用户能够更灵活地指定舞蹈的风格、情绪基调与动作偏好，进一步拓宽三维数字人舞蹈生成的创意表达空间。同时，面向虚拟直播与实时社交等交互式场景，从模型轻量化与高效渲染等角度提升推理效率，推动三维数字人舞蹈生成技术从离线生产向实时交互的应用转型，也是未来值得持续关注的重要课题。